\section{Introduction}

Large language models are often aligned to refuse harmful requests while still
answering benign requests. Refusal ablation is an intervention that attempts to
remove or weaken this refusal behavior. This makes it a useful case study for
statistical evaluation because the intervention may create a safety--utility
trade-off: fewer refusals can improve helpfulness on benign prompts, but can also
increase compliance with unsafe prompts.

This report evaluates whether a refusal-ablated model behaves differently from
the untouched model when both are tested on the same prompt set. The focus is
not on training a new model, but on designing and analyzing a paired model
evaluation with appropriate metrics and statistical tests.

\subsection{Research Questions}

We consider three research questions:

\begin{itemize}
  \item \textbf{RQ1:} Does refusal ablation reduce refusals on unsafe prompts?
  \item \textbf{RQ2:} Does refusal ablation reduce false refusals on benign prompts?
  \item \textbf{RQ3:} Does the ablation effect vary across prompt categories?
\end{itemize}

\subsection{Contribution}

The contribution of the report is a statistical comparison of the untouched and
refusal-ablated model using paired binary response labels. The same prompts are
evaluated on both models, which allows direct comparison of model behavior at
the prompt level. Following the model-comparison setting described by
\citet{raschka2018model}, we use McNemar's test for paired binary outcomes and
report effect sizes alongside significance tests.
